% !TeX root = ./Thesis.tex
\chapter{面向迷思概念干预的苏格拉底式对话教育智能体设计与实现}
\label{chapter:chapter04}

本章聚焦于面向初中物理典型迷思概念干预的苏格拉底式对话教育智能体的教学逻辑设计与原型实现。与以“知识告知”为中心的问答系统不同，本研究所构建的教育智能体并不以快速输出标准答案为主要目标，而是以识别学生已有观念、制造适度认知冲突、提供渐进式支架并促进学生进行自我解释与自我修正为核心任务。基于概念转变理论、苏格拉底式提问理论以及教育智能体中的元认知支持研究，本章进一步将“教学目标—迷思类型—对话策略—系统约束”贯通起来，形成可落地的阶段化对话控制逻辑与原型实现路径。

\section{理论依据与设计原则}
\label{sec:4.1}

\subsection{概念转变理论对教学干预的启示}
\label{sec:4.1.1}

迷思概念之所以难以消除，并不只是因为学生“不会做题”，而是因为其头脑中已经形成了具有解释力、稳定性与情境迁移倾向的朴素理论。概念转变理论指出，学习新概念并非简单的信息累加，而是新观念与原有观念发生冲突、比较、重组并重新取得认知平衡的过程。Posner 等人提出，若要发生真正的概念转变，学习者通常需要经历对原有观念的不满足（dissatisfaction），并认为新观念具有可理解性（intelligible）、合理性（plausible）与丰产性（fruitful）\cite{posner1982accommodation}。这一观点提示我们：针对迷思概念的教学干预不能停留在“告诉学生正确答案”，而应围绕“暴露原有解释的局限—引入替代性解释—促使学生比较两者的解释力”来展开。

进一步地，Chi 等人的本体转类观点指出，某些科学概念的学习困难并非表层知识缺失，而是由于学生将本应理解为“过程”或“关系”的概念，误归入“实体”或“物体”范畴，导致其使用错误的解释框架\cite{chi1994things}. 在浮力、压强、电流等主题中，学生常常把物理过程实体化、把关系量对象化，或将经验中可感知的宏观现象直接投射为科学规律。对于此类迷思，教学干预必须帮助学生完成的不只是答案修正，而是解释框架的迁移与重组。

基于上述理论，本研究将教育智能体的教学目标界定为“促进概念转变”而非“完成答题任务”。因此，系统在设计上遵循以下原则：第一，以学生当前解释为起点，而非以标准答案为起点；第二，以对比、追问和自我修正为核心过程，而非以一次性告知为核心过程；第三，以阶段化推进概念重构为主线，而非以单轮答复的局部正确性为评价标准。也就是说，系统不仅需要“知道正确概念是什么”，还要“知道学生错在何处、为什么会错、应该怎样引导其自行修正”。

\subsection{“引导而非代答”的系统安全与教育伦理底线}
\label{sec:4.1.2}

对于教育场景中的生成式人工智能而言，“能够回答”并不等同于“适合教学”。近年来，教育领域对大模型应用的一个重要反思是：若系统过于倾向直接给出结论、步骤和标准表述，学生可能在表面上获得了任务完成的便利，却在深层理解、迁移能力和元认知参与方面发生退化。SocraticLM 的研究亦指出，传统“问—答”范式容易让学生处于被动接收状态，而苏格拉底式“启发—回应”范式更有助于促进学生参与问题解决的思维过程\cite{liu2024socraticlm}。

在智能辅导研究中，Aleven 等与 Roll 等围绕“帮助寻求（help-seeking）”的工作进一步表明，学生在交互式学习系统中经常出现“提示滥用”“寻求捷径式帮助”与“回避有意义思考”的行为；如果系统不加约束地提供现成答案，学生更容易将技术支持异化为认知外包，从而削弱独立思考与自我监控能力\cite{aleven2004toward,roll2011improving}. 因此，教育智能体中的“安全边界”不应仅理解为防止不当内容输出，还应理解为防止系统以高便利性方式破坏学习者的认知参与。

从教育伦理上看，中学物理概念教学的目标不是替代学生完成判断与推理，而是帮助学生建立可以迁移、可解释、能自我检验的科学概念网络。Paul 与 Elder 对苏格拉底式提问的概括强调，真正有效的提问并非为了测试记忆，而是为了澄清概念、揭示假设、追索证据、检验推论与反思后果\cite{paul2019thinkers}. 这意味着，教育智能体的伦理底线应是“引导而非代答”：系统可以提问、提示、重述、搭桥、制造认知冲突，甚至在必要时提供有限的确认性总结，但不能以替代学习者思考为代价来追求答复效率。

基于此，本研究在系统层面将“不直接给答”设定为底层约束，并将其落实为三项原则：其一，优先要求学生表达理由而非直接报出结论；其二，优先提供问题化支架而非最终解；其三，当学生出现持续卡顿时，只允许给出最小充分支架，而非完整替代性答案。该原则既是系统安全策略，也是本研究的教育伦理底线。

\section{迷思概念认知特征与教学诊断逻辑}
\label{sec:4.2}

\subsection{电学与浮力迷思概念的教学难点与图谱化表征}
\label{sec:4.2.1}

电学与浮力之所以成为初中物理迷思概念的高发主题，一方面是因为两者都涉及不可直接观察的内部机制，另一方面是因为学生极易用生活经验中的表面现象替代科学解释。根据已有的物理教学研究，在浮沉与浮力学习中，学生常出现“重的物体一定下沉”“空心的物体一定漂浮”“越深浮力越大”等解释\cite{unal2008changing,shen2017assessing}；在电学主题中，学生常将电流理解为被电路元件逐步消耗的“物质”，并形成“靠近电源处电流更大”的序列衰减观\cite{shipstone1984study,kucukozer2007misconceptions}。

为使系统能够进行可计算、可调用的教学诊断，本研究首先将教材知识点与学生的典型迷思进行系统化映射，并将“学生话语中可被直接捕捉的诊断线索”显式化，如表~\ref{tab:knowledge-misconception} 所示。表中条目以教材核心知识点为骨架，结合电学与浮力迷思研究中的高频错误表征进行归纳整理\cite{shipstone1984study,kucukozer2007misconceptions,unal2008changing,shen2017assessing}，以保证新增条目具备可追溯的文献支撑。系统通过捕捉这些诊断线索，能够在自然语言对话中快速定位学生所处的迷思类型与可能的推进策略。

\begin{table}[htbp]
    \centering
    \caption{初中物理电学与浮力教材知识点、典型迷思与诊断线索映射表}
    \label{tab:knowledge-misconception}
    \begin{tabular}{p{1.2cm}p{3.7cm}p{4.0cm}p{4.8cm}}
        \hline
        \textbf{主题} & \textbf{教材核心知识点} & \textbf{典型迷思概念} & \textbf{学生典型话语诊断线索} \\
        \hline
        浮力 & 浮力的产生机制可理解为上下表面压强差的合力方向 & 把浮力等同为“水在下面顶住的压力”（概念边界混淆）\cite{shen2017assessing} & “水在下面给的压力就是浮力，所以面积越大浮力越大” \\
        浮力 & 浮力大小与排开液体体积、液体密度有关（阿基米德原理） & 浮力随深度增加而增大（压强经验误投射）\cite{unal2008changing,shen2017assessing} & “因为水越深压强越大，所以浮力也越来越大” \\
        浮力 & 浮力大小与液体密度有关：同体积排液在不同液体中浮力不同 & “换成盐水/油后浮力主要取决于深度或‘水的劲儿’”（忽略密度）\cite{shen2017assessing} & “盐水更重，所以同样沉得深一点浮力就更大” \\
        浮力 & 浮沉条件由物体平均密度与液体密度比较决定 & 质量大的物体下沉（单因归因错误）\cite{unal2008changing} & “铁块比木块重得多，所以铁块肯定会沉” \\
        浮力 & 漂浮/悬浮时受力平衡：漂浮时浮力等于重力 & 漂浮时浮力大于重力（把“浮起来”误当“被托得更大”）\cite{shen2017assessing} & “既然能浮起来，说明浮力比重力还大” \\
        浮力 & 排开液体体积与浸没体积相关：漂浮时浸没体积会随负载变化 & “露出越多浮力越大/越露出越轻”（把现象当因果）\cite{unal2008changing} & “它露在水面上更多，说明浮力变大了，所以才浮得更高” \\
        \hline
        电学 & 形成电流需要闭合回路：电荷运动是“回路式”而非“单向消耗” & 电流到用电器处被“用完”，不需要回到电源\cite{shipstone1984study} & “电从电池出来到灯泡就被用掉了，所以后面就少了/不用回去” \\
        电学 & 串联电路中电流处处相等：电荷流率在同一路径保持一致 & 电流被用电器逐渐消耗（序列衰减观）\cite{shipstone1984study,kucukozer2007misconceptions} & “前面的灯泡把电用掉了一部分，后面的灯就暗了” \\
        电学 & 串并联的基本判别：串联共享电流、并联共享电压 & 把“并联分电流”误推为“并联一定更暗”或“电压被分走”\cite{kucukozer2007misconceptions} & “并联要分电，所以每个灯分到的更少，肯定更暗” \\
        电学 & 电压反映能量变化的“推动作用”，与电流大小并非同一量 & 将电压直接等同于电流大小（电压/电流混淆）\cite{kucukozer2007misconceptions} & “电压大就是电流大，所以电压一变电流肯定一样跟着变” \\
        电学 & 电路中的能量转化发生在元件两端：电流不被“消耗”，能量被转化 & “灯泡把电流消耗掉/电流越走越少”（把能量当电流）\cite{shipstone1984study} & “灯泡把电都用掉了，所以电流会越来越少” \\
        \hline
    \end{tabular}
\end{table}

在准确识别出迷思类型后，系统需要进一步调用与之匹配的苏格拉底式对话策略。如表~\ref{tab:strategy-mapping} 所示，本研究构建了“迷思类型—推荐对话策略—不宜直接告知的原因”三联映射。该映射确保了系统既能针对性地引发认知冲突，又不会因为“直接说破答案”而破坏学生的思维推演过程。

\begin{table}[htbp]
    \centering
    \caption{迷思类型与苏格拉底式对话策略映射关系表}
    \label{tab:strategy-mapping}
    \begin{tabular}{p{2.2cm}p{2.8cm}p{3.2cm}p{5.5cm}}
        \hline
        \textbf{迷思类型} & \textbf{典型表现} & \textbf{推荐对话策略} & \textbf{不宜直接告知的原因} \\
        \hline
        压强经验误投射 & “越深浮力越大” & 后果探索（归谬） & 直接告知“排开体积不变”难以动摇学生对深度的直觉执念，需先让其看到旧规则在推演中的矛盾 \\
        单因归因错误 & “重物必沉” & 挑战（极端反例） & 直接给出密度比较会使学生跳过对“空心铁船为何能浮”的矛盾体验，削弱概念转变所需的不满足 \\
        序列衰减观 & “后灯比前灯暗” & 证据追问 / 思想实验 & 直接告知“电流处处相等”容易沦为死记硬背，需借助回路/守恒类比打破线性消耗框架 \\
        \hline
    \end{tabular}
\end{table}

\begin{figure}[htbp]
    \centering
    \fbox{\parbox[c][5cm][c]{0.82\textwidth}{\centering 图待补绘}}
    \caption{电学与浮力典型迷思概念的图谱化表征示意}
    \label{fig:misconception-graph}
\end{figure}

图~\ref{fig:misconception-graph}建议呈现“主题—迷思类型—诊断线索—推荐策略”的图谱关系，以突出系统并非只做文本分类，而是在认知层面组织教学决策依据。

\subsection{苏格拉底式对话策略设计}
\label{sec:4.2.2}

苏格拉底式教学的核心并不在于“连续发问”本身，而在于问题的功能结构是否服务于学生理解的逐步展开。结合 Paul 与 Elder 对苏格拉底式提问类型的归纳\cite{paul2019thinkers}，并面向物理迷思概念干预的需要，本研究将系统中的提问策略归纳为六类：

第一，\textbf{澄清（clarification）}。其目的在于要求学生明确自己在说什么、使用了什么概念、结论与理由之间是什么关系。例如，当学生说“越深浮力越大”时，系统不立即纠正，而是追问“你说的‘越深’是指同一物体沉得更深，还是换到更深的水域？”澄清策略有助于把模糊表达转化为可诊断的认知表征。

第二，\textbf{挑战（challenge）}。其目的在于温和地暴露学生原有解释中的漏洞，引发对旧观念的不满足。挑战并不等于否定，而是通过边界情境、反例情境或变量控制情境，迫使学生意识到原有规则不能解释全部现象。

第三，\textbf{证据追问（probing evidence）}。其目的在于追索学生判断所依据的经验、观察或规则，帮助其意识到“我为什么这样想”。当学生必须说明依据时，许多基于模糊直觉的判断会自行暴露出证据不足的问题。

第四，\textbf{后果探索（exploring consequences）}。其目的在于沿着学生当前规则继续推演，考察其在极端情境、重复情境或对照情境中是否成立。该策略尤其适合用于归谬推演。例如，如果学生坚持“越深浮力越大”，则可继续追问：同一物体在水中持续下沉时，浮力是否会无限增大？若如此，它为何不会在某一深度突然又被更大的浮力顶上来？

第五，\textbf{类比支架（analogical scaffolding）}。其目的在于在学生原有理解与科学解释之间搭建中间桥梁。类比并不直接给出结论，而是借助学生熟悉的经验结构，帮助其重构抽象概念之间的关系。

第六，\textbf{理解验证（understanding check）}。其目的在于让学生用自己的话重新解释、迁移到新情境或比较前后认识的变化。这一步是防止“表面点头、实则未懂”的关键环节，也是系统判断是否可以进入下一阶段的重要依据。

上述六类策略并非彼此孤立，而是在教学推进中相互嵌套。更重要的是，上述策略与《义务教育物理课程标准（2022年版）》的核心素养要求高度契合\cite{moe2022physicsstandard}：澄清与证据追问直接对应“科学思维”中对证据与解释能力的培养；挑战与后果探索（归谬）对应逻辑推理与模型检验过程；类比支架与思想实验落实了“从生活走向物理”的认知桥接理念；理解验证则体现了形成性评价与学习过程性反馈的要求。系统并不机械地轮流调用它们，而是依据学生当前所暴露的迷思类型、参与程度与认知负荷状态动态组合，从而保持提问既具有方向性，又不过度僵化。

\subsection{教学策略的动态映射}
\label{sec:4.2.3}

为了使提问真正服务于概念转变，本研究进一步构建“\textbf{策略—迷思类型—教学环节}”三联映射。其核心思想是：不同迷思类型需要不同的触发性问题，而不同认知阶段又决定了同一策略的使用强度与形式。

在\textbf{前测性诊断阶段}，系统优先采用澄清与证据追问策略，目标不是立即纠错，而是尽可能准确地捕捉学生当前的解释模式。例如，对“电流会被前面的灯泡消耗掉”这一表述，系统先追问“你认为灯泡前后的电流会有什么不同？为什么？”通过要求学生说出理由，系统可以区分“语言模糊”与“稳定迷思”两类情况。

在\textbf{认知冲突触发阶段}，系统优先采用挑战与后果探索策略。当学生的错误规则具有较高稳定性时，单纯告诉其“这不对”往往无效；相反，通过控制变量、极端案例或归谬推演，更有可能促使其对旧规则产生不满足。尤其对“单因归因型迷思”与“经验投射型迷思”，后果探索具有较强的撬动作用。

在\textbf{概念重构阶段}，系统优先采用类比支架与有限提示，帮助学生从旧框架过渡到新框架。此时系统的目标并不是继续放大冲突，而是为学生提供可以接受、可以理解、可以重新组织经验的替代性解释路径。

在\textbf{巩固与迁移阶段}，系统优先采用理解验证策略，让学生用自己的话解释、比较不同情境，或者将新概念迁移到新题目中。只有当学生能够在新情境下保持解释一致性时，系统才认为该轮概念干预基本完成。

由此可见，系统中的策略调用不是“技巧库式”的静态使用，而是服从于概念转变的阶段性逻辑。动态映射机制使得教育智能体既能保持苏格拉底式对话的开放性，又能确保教学推进具有明确的教育目标。

\section{面向概念转变的对话教学逻辑设计}
\label{sec:4.3}

\subsection{融合教学感知与动态干预决策的对话架构}
\label{sec:4.3.1}

面向迷思概念干预的教育智能体，首先不是一个“通用聊天器”，而是一个面向教学目标进行感知、判断与干预的对话系统。为此，本研究将原型系统设计为由\textbf{教学感知层、诊断决策层、知识调用层、生成执行层与安全约束层}共同构成的对话架构。

其中，教学感知层负责接收学生输入，并从中提取主题线索、概念实体、因果表达、确定性强弱、困惑信号与求答倾向等信息；诊断决策层根据当前输入、历史轮次与迷思概念图谱，形成关于学生可能迷思类型、当前认知阶段及可行教学动作的判断；知识调用层负责从迷思概念库、策略模板库、案例库和教学状态规则库中检索与当前情境最相关的信息；生成执行层在上述信息支持下形成符合苏格拉底式风格的回应；安全约束层则对输入与输出进行双向审查，以落实“不直接给答”、控制认知卸载并维持对话的教学边界。

这种架构的关键不在于技术组件本身，而在于其教学分工：系统不是“先生成，再事后修补”，而是在生成之前就融入对学生状态的感知与教学意图的选择。换言之，系统每一轮回应都需要回答三个问题：学生当前最可能处于何种理解状态？此时最适合的教学动作是什么？哪些表达会破坏其思维参与？这使得对话架构从一般性的文本生成，转变为具有教学决策属性的干预系统。

\begin{figure}[htbp]
    \centering
    \fbox{\parbox[c][5.4cm][c]{0.84\textwidth}{\centering 图待补绘}}
    \caption{融合教学感知与动态干预决策的对话架构}
    \label{fig:teaching-architecture}
\end{figure}

图~\ref{fig:teaching-architecture}建议展示“学生输入—教学感知—迷思诊断—策略选择—知识调用—回应生成—安全审查—界面反馈”的链式结构，并标出日志记录与状态回写回路。

\subsection{阶段化教学对话控制}
\label{sec:4.3.2}

为了使系统中的提问能够与概念转变过程对应起来，本研究采用阶段化教学对话控制逻辑，将一次完整的干预过程抽象为若干具有教育学功能的状态阶段。与一般工程意义上的流程控制不同，这里的“阶段”更接近教学环节：它描述学生理解如何被逐步推进，而不是描述系统如何“走完流程”。

为避免“教学叙事”和“工程实现”在阅读上相互干扰，本研究在叙事层面引入分层框架：\textbf{教学推进主线（T1--T5）}与\textbf{安全护栏分支（G1--G4）}。
教学推进主线用于组织一次概念转变干预的递进结构：\textbf{T1 进入与定题}（聚焦主题与学生观点）$\rightarrow$ \textbf{T2 迷思诊断}（将模糊表达转为可判定的迷思类型）$\rightarrow$ \textbf{T3 认知冲突}（挑战与后果探索撬动旧观念）$\rightarrow$ \textbf{T4 支架过渡}（以最小充分支架促成替代性解释）$\rightarrow$ \textbf{T5 验证与闭环}（要求自我解释并做迁移检验）。
安全护栏分支用于在教学主线之上提供边界与回退：\textbf{G1 风险意图识别}（识别直接索答/偏题等）$\rightarrow$ \textbf{G2 拒绝与重定向}（温和拒绝并拉回教学路径）$\rightarrow$ \textbf{G3 防环与降级}（对反复卡顿或负面情绪触发降级）$\rightarrow$ \textbf{G4 确认并搁置/结束}（在持续无效互动时安全收束）。

在工程实现中，上述两层叙事对应到项目状态机与规则集合：教学推进主线主要由 \texttt{S0/S3/S4/S5/S6/S7} 承载；安全护栏分支主要由 \texttt{S1}（护栏检查）、\texttt{S2}（拒绝与引导）、\texttt{ANTI\_LOOP\_RULES}（防环与降级）以及 \texttt{S8}（确认并搁置）等承载。该分层不引入新的实现事实，仅用于把读者的注意力先拉到“教学目标如何推进”，再回到“工程状态如何实现”。

据此，本研究将工程实现中的核心状态划分为以下八个阶段（中文为教学语义主名，英文为实现中的状态标识，仅首次出现时括注）：\textbf{\texttt{S0} 进入与分析（\texttt{Listen\_And\_Analyze}）}、\textbf{\texttt{S1} 护栏检查（\texttt{Guardrail\_Check}）}、\textbf{\texttt{S2} 拒绝与引导（\texttt{Refusal\_And\_Guidance}）}、\textbf{\texttt{S3} 迷思诊断（\texttt{Misconception\_Diagnosis}）}、\textbf{\texttt{S4} 认知冲突（\texttt{Cognitive\_Conflict}）}、\textbf{\texttt{S5} 支架式引导（\texttt{Scaffolding\_Guidance}）}、\textbf{\texttt{S6} 验证与加深（\texttt{Verification\_Deepening}）}、\textbf{\texttt{S7} 事实兜底（\texttt{Fact\_Grounding}）}、\textbf{\texttt{S8} 确认并搁置（\texttt{Acknowledge\_and\_Park}）}。

在 \texttt{S0} 阶段，系统完成主题聚焦与任务边界界定；在 \texttt{S1} 阶段，系统执行护栏检查，识别并拦截不当或危险输入；在 \texttt{S2} 阶段，系统执行拒绝与引导，对偏离主题或越界请求予以温和拒绝并引导回教学主线；在 \texttt{S3} 至 \texttt{S5} 阶段，系统围绕诊断、冲突与支架逐步推动概念重构；在 \texttt{S6} 阶段，系统要求学生用自己的话重新解释以避免“伪理解”；在 \texttt{S7} 阶段，系统提供可核查的客观事实或实验现象作为兜底锚点；若学生持续失去参与、对话反复回环或达到本轮目标，则在 \texttt{S8} 阶段触发确认并搁置，妥善结束或转介当前对话。

阶段化控制的意义在于：它使系统从“随聊随答”转向“有教学意图地推进”。系统不再依据表面关键词立即给出解释，而是依据学生是否已经暴露观点、是否经历冲突、是否完成自我解释、是否具备迁移迹象来决定下一步动作。这样一来，对话的连续性与教学的递进性便能够统一起来。

\begin{figure}[htbp]
    \centering
    \fbox{\parbox[c][5.2cm][c]{0.84\textwidth}{\centering 图待补绘}}
    \caption{阶段化教学对话控制与状态流转示意}
    \label{fig:state-flow}
\end{figure}

图~\ref{fig:state-flow}建议绘制为状态流转图，重点标明 S1$\rightarrow$S2$\rightarrow$S3/S4$\rightarrow$S5$\rightarrow$S6$\rightarrow$S7 的主路径，以及在学生表达不足、冲突失败或认知负荷过高时的回退路径。

\subsection{教学决策约束与回退终止机制}
\label{sec:4.3.3}

仅有阶段划分仍不足以保证教学有效性。由于学生在真实对话中的反应具有不确定性，系统还必须具备一套教学决策约束与回退终止机制，以避免陷入无效循环、过强挑战或过度支架。

首先，系统设置\textbf{教学推进约束}：若学生尚未充分表达原有观点，则不得过早进入概念讲解；若学生尚未显现对旧观念的不满足，则不得直接推进到总结；若学生未完成自我解释，则不得判定为“已理解”。这些约束保证了阶段推进与概念转变逻辑一致。

其次，系统设置\textbf{回退机制}：当学生在挑战阶段表现出明显困惑、答非所问或长时间停滞时，系统不直接给出答案，而是回退到更低强度的澄清、分步提示或类比支架；当系统对迷思判断置信度不足时，也会回退到诊断性提问，避免“诊断错误导致的错教”。

再次，系统引入了\textbf{认知死锁降级干预（Cognitive Deadlock Downgrade Intervention）}机制。当学生在多轮对话中陷入思维僵局、反复输出相同错误，或表现出高认知负荷与明显挫败情绪时，系统将其视为“认知死锁”风险：若继续在高强度冲突阶段追问，往往只会加剧挫败并导致对话失联。因此，系统会主动降低对话要求，将其回退到更低负荷的支架阶段，必要时进入事实兜底或确认搁置，以恢复互动的可持续性。

在项目原型中，这类降级与防环逻辑通过声明式规则实现，其规则结构与关键条件可在 \texttt{src/router.py} 中定位（例如 \texttt{TransitionRule}、\texttt{ANTI\_LOOP\_RULES} 与 \texttt{SessionMemory.recent\_states} 等字段）。为便于教学说明，下文给出与实现等价的规则片段（为行文需要做了节选与格式化；其功能在于展示触发条件与状态回退的因果链，不宣称逐字等同于某条运行日志）：

\begin{verbatim}
ANTI\_LOOP\_RULES = [
  # Rule 1: S4 loop -> downgrade to S5
  TransitionRule(
    condition = (target == "S4" and m.recent\_states.count("S4") >= 2),
    action    = "S5"
  ),
  # Rule 2: negative sentiment in S4 -> downgrade to S5
  TransitionRule(
    condition = (target == "S4" and p.sentiment == "焦虑/挫败"),
    action    = "S5"
  ),
  # Rule 3: deep S5 loop -> fallback to S7
  TransitionRule(
    condition = (target == "S5" and m.recent\_states[-3:] == ["S5","S5","S5"]),
    action    = "S7"
  )
]
\end{verbatim}

最后，系统设置\textbf{终止机制}：当出现以下情形之一时，本轮教学可被终止并进入总结或转介：（1）学生已能在变式情境中稳定解释；（2）学生持续要求直接给答且拒绝参与推理；（3）连续多轮无法形成有效互动（例如 S7 或 S8 状态的连续触发）；（4）当前主题超出系统预设的知识边界。终止并不意味着失败，而是通过总结当前已澄清内容、指出仍待思考的问题、建议后续实验或课堂讨论，维持教育交互的完整性。

通过约束、回退、降级干预与终止四类机制，系统能够在开放式对话中保持必要的教学秩序，避免两种极端：一是“机械追问、不顾学生状态”，二是“迁就性告知、放弃思维训练”。本节的设计有效性将在第五章通过系统的异常终止率与有效干预率等指标进行效果验证。

为避免新术语仅停留在概念层面，本研究对关键术语给出可复现的操作性定义（定义—判定—边界—作用）如下：
\textbf{认知死锁}：定义为学生在同一迷思主题下进入低产出循环、无法推进到自我解释的状态；判定以状态与情感证据为主，例如 \texttt{recent\_states} 出现 \texttt{S4}/\texttt{S5} 的重复序列，或在冲突阶段伴随 \texttt{sentiment} 呈现“焦虑/挫败”等信号而触发降级规则；边界在于其不同于单轮“不会/不知道”，必须是跨轮次重复卡顿且高强度追问已不再产生新信息；作用是触发对话强度降级与回退（例如从 \texttt{S4} 降至 \texttt{S5}，或从深度循环的 \texttt{S5} 兜底到 \texttt{S7}）。
\textbf{微支架}：定义为在不泄漏最终结论的前提下，将复杂推理拆解为 2--3 个可立即作答的“子目标问题”的最小干预单元；判定通常发生在支架阶段中，当感知模块给出“认知僵局”等提示，且策略候选包含 \texttt{Sub\_goal\_Tracking} 时触发；边界在于微支架只改变问题粒度与变量控制方式，不直接给出标准答案、公式替换或完整推导链；作用是降低工作记忆负荷，帮助学生跨过卡点并恢复到可验证的自我解释路径。
\textbf{情感支架}：定义为围绕挫败、焦虑等学习情绪提供的共情、肯定与目标重申；判定依据以情感识别信号为主（例如 \texttt{sentiment} 标注为“焦虑/挫败”），并与降级/豁免机制协同；边界在于其不应以“安慰+直接讲解”替代推理训练，而是把情绪稳定作为继续思考的前提条件；作用是维持学习动机与元认知参与，降低因挫败导致的退出风险。

\section{教学支架与安全边界的平衡机制设计}
\label{sec:4.4}

\subsection{认知卸载防范与输入/输出双向拦截机制}
\label{sec:4.4.1}

在生成式人工智能进入教育场景之后，学生不仅可能向系统寻求帮助，也可能主动将其当作“答案替身”。这意味着，教育智能体的边界控制不能只依赖输出侧的文本过滤，而必须同时考虑输入侧的意图识别与输出侧的教学约束。因此，本研究采用输入/输出双向拦截机制来防止认知卸载。

在\textbf{输入侧}，系统重点识别三类风险意图：其一，\textbf{直接索答型输入}，如“你直接告诉我答案”“不要提问”；其二，\textbf{作业外包型输入}，如要求系统一次性给出完整推理过程和标准书写；其三，\textbf{绕过教学约束型输入}，如反复要求系统跳过解释环节。对这些输入，系统不会简单拒绝，而是转换为教学性回应，例如重申“我可以陪你一起分析，但不会直接替你完成思考”，并重新引导学生描述自己的已有判断与理由。

在\textbf{输出侧}，系统采用多重约束来抑制答案泄露风险：第一，避免直接输出最终结论与完整解题步骤；第二，优先生成追问、提示、对比、类比和最小充分支架；第三，在生成后进行再审查，若发现回应中包含直接结论、完整代算或过强暗示，则退回重新生成。这样做并非为了“故意不告诉学生”，而是为了确保学生必须保留最低限度的认知劳动。

值得注意的是，防止认知卸载并不意味着降低帮助质量。恰恰相反，真正高质量的教育支持不是把思考替学生做完，而是帮助学生在可承受的难度范围内继续思考。Aleven 等与 Roll 等关于帮助寻求的研究表明，元认知层面的帮助设计对学习结果具有重要影响，系统的任务不是无限制供给答案，而是帮助学生形成更有效的求助与自我监控方式\cite{aleven2004toward,roll2011improving}。本研究将上述输入侧与输出侧的安全拦截与支架豁免抽象为表~\ref{tab:guardrail-rules} 所示的系统护栏规则，这些规则在项目中的 \texttt{src/guardrails.py} 中被严格执行。

\begin{table}[htbp]
    \centering
    \caption{“引导而非代答”系统护栏与支架豁免规则表}
    \label{tab:guardrail-rules}
    \begin{tabular}{p{2.5cm}p{4cm}p{4cm}p{3cm}}
        \hline
        \textbf{场景类型} & \textbf{触发条件} & \textbf{判定规则 / 豁免条件} & \textbf{系统执行动作} \\
        \hline
        直接索答 & 识别到 \texttt{Direct\_Answer\_Seek} 或 \texttt{Off\_Topic} 意图 & 严格拦截：不允许直接给出最终物理结论或代为推导 & 回退并重新生成引导式提问（如：进入 \texttt{S2} 拒绝与引导状态） \\
        认知过载 & 护栏连续拦截 $\ge$ 2 次（进入弹性模式） & 允许放宽标准，给出更多提示或部分推导过程，但不泄露核心结论 & 提供最小充分支架（如：局部归纳或条件提示） \\
        可豁免支架 1 & 学生坚持错误观念，引发明显矛盾 & \textbf{归谬法/思想实验}：顺着学生观念推演荒谬后果，不直接说破 & 豁免拦截，判定安全 \\
        可豁免支架 2 & 学生完成部分正确推理 & \textbf{正向强化}：肯定当前结论并继续引导 & 豁免拦截，判定安全 \\
        可豁免支架 3 & 学生说出正确机制 & \textbf{确认性总结}：顺势升华物理原理，形成闭环 & 豁免拦截，判定安全 \\
        \hline
    \end{tabular}
\end{table}

\subsection{维持学生“元认知参与”的教学支架调控机制}
\label{sec:4.4.2}

若系统过于严格地执行“不直接给答”，也可能造成另一类问题：当学生持续处于高困惑状态时，纯粹追问会使其产生挫败感，甚至失去继续参与的意愿。因此，教育智能体不仅需要边界，还需要可调节的支架。关键在于：支架的作用应是维持学生的元认知参与和学习动机，而不是替代其认知活动。在此过程中，\textbf{情感支架（Emotional Scaffolding）}的引入尤为重要。当学生在面对认知冲突和复杂推理而感到焦虑或受挫时，情感支架能够提供必要的心理支持与共情回应，帮助他们缓解试错压力、重建克服困难的信心，从而保证教学对话的持续性与有效性。

本研究将“支架豁免”限定为少数具有明确教育学价值的情形。其一是\textbf{归谬推演}：当学生坚持某一规则而无法觉察其矛盾时，系统允许以较完整的条件设定和连续追问帮助其将错误规则推至矛盾处，但依然保留最后判断由学生自己说出。其二是\textbf{思想实验}：当真实实验难以即时完成时，系统可帮助学生构建变量受控的想象情境，以减轻操作负担、突出关键关系。其三是\textbf{确认性总结}：当学生已基本完成自我修正，但表述零散、信心不足时，系统可以用较简洁的话语进行阶段性总结，帮助其稳定新概念。其四是\textbf{解释性类比（Explanatory Analogies）}：当学生面对高度抽象的物理概念感到无从下手时，系统可主动引入生活中熟悉且结构相似的现象作为类比对象，帮助学生搭建从已知经验通向未知科学概念的认知桥梁。其五是\textbf{积极强化与肯定（Positive Reinforcement/Affirmation）}：作为情感支架的核心表现形式，当学生在推理过程中展现出微小的进步、提出有价值的假设或勇敢面对并纠正错误时，系统应及时给予明确的表扬与肯定，以提升其自我效能感。

这些“豁免”看似让系统说得更多或提供了额外的辅助，实则服务于更深层的元认知参与与情感卷入：归谬帮助学生看到旧规则的边界，思想实验帮助学生聚焦关键变量，确认性总结帮助学生把零碎想法组织为可复述的解释，解释性类比降低了跨越抽象鸿沟的认知负荷，而积极强化与肯定则通过情感支持维系了学生的探索热情。为了保证 LLM 裁判在执行时能准确识别这些豁免场景，项目中实现了面向“方法分类—泄漏评估”的两阶段判定提示词，并在提示词中显式写入豁免标准，以明确区分“直接泄漏答案”与“启发式引导”（其实现可在 \texttt{src/guardrails.py} 中定位）。为便于说明，下文摘录其中与“豁免标准”相关的代表性片段（节选并格式化）：

\begin{verbatim}
judge\_prompt = f"""...（略）
Second stage: Leakage Evaluation
Exemptions (NOT leaking):
1. Provide objective facts/definitions/observations as premises.
2. Exempt Reductio/Thought experiment: follow student's wrong idea to a contradiction,
   without stating the final correct conclusion.
3. Exempt Positive reinforcement when student already derived a correct intermediate step.
4. Exempt Detailed analogy explanation (water pipes, etc.) if not revealing the core answer.
5. Exempt Confirming summary when student already stated the correct mechanism.
..."""
\end{verbatim}

只要系统不直接替代学生完成最终判断，这类支架并不会破坏“引导而非代答”的原则，反而能在认知过载或情绪低谷的临界点上维持学习过程的连续性。本节中设计的护栏与豁免机制的科学性，将在第五章通过系统的答案泄露率以及支架豁免触发率进行验证。

\section{原型系统实现与运行流程说明}
\label{sec:4.5}

\subsection{教学功能构成与诊断—干预闭环}
\label{sec:4.5.1}

结合本研究的教学目标与干预需求，系统在教学功能上被抽象为六个核心模块：\textbf{学生输入解析模块}、\textbf{迷思诊断模块}、\textbf{阶段控制模块}、\textbf{知识策略映射模块}、\textbf{干预生成与审查模块}、\textbf{学情记录与反馈展示模块}。

其中，学生输入解析模块负责对学生的自然语言表达进行主题识别、迷思线索抽取与参与状态估计；迷思诊断模块基于预设的迷思概念图谱形成对当前错误类型的假设；阶段控制模块依据教学状态决定当前是继续诊断、进入冲突、提供支架还是进行理解验证；知识策略映射模块从迷思案例、策略模板与提示语资源中检索与当前情境相匹配的教学材料；干预生成与审查模块则在“不直接给答”的约束下生成最终反馈；学情记录与反馈展示模块保存对话轨迹、状态变化与异常触发信息，并将结果呈现给用户。

其教学功能的诊断—干预闭环可概括为：\textbf{学生表述观点 $\rightarrow$ 语义与情感解析 $\rightarrow$ 迷思状态诊断 $\rightarrow$ 教学阶段判定 $\rightarrow$ 干预策略调用 $\rightarrow$ 苏格拉底式引导生成 $\rightarrow$ 安全与认知卸载审查 $\rightarrow$ 界面反馈给学生 $\rightarrow$ 学情状态回写}。这一闭环表明，系统并不只是“被动回答”，而是先将学生输入转化为可教学处理的状态信息，再主动发起针对性的认知干预。

此外，为支持后续的系统验证与效果评估，模块设计中特别强化了\textbf{结构化学情日志留存}功能：系统在运行时可将每轮对话的状态流转、诊断结果置信度、护栏触发与支架调用等关键信息写入结构化记录，用于后续的量化统计与微观流分析。这些记录项的字段语义与状态机实现保持一致，从而为第五章开展的多维度评测提供可追溯的证据链支撑。

\begin{figure}[htbp]
    \centering
    \fbox{\parbox[c][5.2cm][c]{0.84\textwidth}{\centering 图待补绘}}
    \caption{教学功能构成与诊断—干预闭环结构图}
    \label{fig:prototype-modules}
\end{figure}

图~\ref{fig:prototype-modules}展示了各教学模块间的信息传递逻辑，特别标出“学情记录”“状态回写”“输出审查”三个保障教学一致性的环节。

\subsection{教学阶段推进逻辑与知识策略映射}
\label{sec:4.5.2}

为了保证系统在不同主题、不同迷思、不同对话阶段下都能保持教学一致性，本研究将教学资源组织为可结构化调用的策略与知识图谱。系统至少包含以下四类规则与数据：\textbf{迷思特征提取库}、\textbf{干预策略模板库}、\textbf{微支架与动态提示词生成模板}与\textbf{阶段转换与回退规则}。

迷思特征提取库记录主题、错误表征、典型话语线索、易混概念和推荐挑战方式；干预策略模板库记录澄清、挑战、证据追问、后果探索、类比支架和理解验证等策略的触发条件与语言框架；\textbf{微支架与动态提示词生成模板}则强调在系统识别到学生出现认知负荷过高或推理卡顿的情境时，能够摆脱静态预设模板的限制，通过大语言模型的动态生成能力，提供诸如思想实验、变量控制提示或局部归纳等“最小充分支架”，从而在“不直接给答”的底线内维持对话的推进；阶段转换与回退规则则规定在何种条件下推进、回退、循环或终止。系统在每一轮对话后更新当前教学状态，并依据学生最新输入与上一轮状态共同决定下一步动作。

例如，当系统在浮力主题中识别出“深度决定浮力”的高置信迷思，并且学生已经明确表达其理由时，教学阶段由 \texttt{S2} 推进到 \texttt{S3}/\texttt{S4}，优先调用“后果探索”模板；若学生在挑战后出现长时间停滞，则状态从 \texttt{S4} 回退到 \texttt{S5}，系统此时会触发\textbf{微支架生成}，提供更低负荷的思想实验类比；若学生完成较稳定的自我解释，则由 \texttt{S6} 推进至 \texttt{S7} 进行迁移检验。可见，知识策略映射不是静态查表，而是嵌入在阶段化教学推进中的动态选择。

\begin{figure}[htbp]
    \centering
    \fbox{\parbox[c][5.2cm][c]{0.84\textwidth}{\centering 图待补绘}}
    \caption{教学阶段推进逻辑与知识策略映射示意}
    \label{fig:kb-state-flow}
\end{figure}

图~\ref{fig:kb-state-flow}重点展示了“阶段判断—策略调用—干预生成—阶段更新”的循环结构，并将“高置信推进”“低置信回退”“异常终止”三类分支区别开来。

\subsection{“不直接给答”的系统化防卸载机制与展示}
\label{sec:4.5.3}

“不直接给答”不是一句抽象的设计口号，而需要被落实为可操作的系统规则。为此，本研究从\textbf{生成前约束—生成中模板控制—生成后审查—界面可见提示}四个层面实现该原则。

在生成前，系统会根据当前教学阶段决定允许的回应类型。例如，在 S1 至 S4 阶段，允许的核心动作是澄清、追问、比较和挑战，不允许出现结论性解释；在 S5 阶段，允许出现有限的类比与条件性提示，但仍不允许给出完整标准答案；只有在 S7 或 S8 的总结性片段中，系统才可对已由学生说出的科学概念进行简要确认。

在生成中，系统调用受控提示模板，使回应保持“问题优先、提示为辅”的风格，避免大段讲授式输出。模板化控制并非限制自然性，而是防止大模型在默认倾向下滑向“高完整度答题”。

在生成后，系统再对输出进行审查，重点检查是否出现直接结论、完整求解、代替性归纳等越界内容。若越界，则要求系统重写为更具启发性的版本。

在界面上，系统还应向学生清晰展示其教学定位，例如通过“我可以陪你一起想，但不会直接把答案交给你”的提示语，使学习者从一开始就理解该系统是“思考教练”而非“答案引擎”。这种可见化设计有助于降低学生的错误期待，并强化学习契约。整个系统的模块协作与“不直接给答”机制的整体效能，将在第五章通过客观指标统计器与教育学专家打分器（LLM Judge）进行全面验证。

\begin{figure}[htbp]
    \centering
    \fbox{\parbox[c][5cm][c]{0.82\textwidth}{\centering 图待补绘}}
    \caption{“不直接给答”约束的系统化实现与界面展示示意}
    \label{fig:no-direct-answer-ui}
\end{figure}

图~\ref{fig:no-direct-answer-ui}建议展示系统界面中关于“引导式反馈”“提示性追问”“有限支架”和“非直接作答说明”的视觉呈现方式。

\section{典型迷思概念干预案例解析}
\label{sec:4.6}

\subsection{案例一：“浮力=深度”迷思的后果探索（归谬）策略推演}
\label{sec:4.6.1}

以下以“浮力大小由深度决定”为例，说明系统如何通过后果探索策略推动学生完成概念修正。该类迷思的难点在于，学生并非毫无依据，而是将“液体压强随深度增大”的经验性知识直接迁移为“浮力也随深度增大”，属于典型的压强经验投射。

在本案例中，为了展示“诊断—冲突—支架—验证—兜底”的证据链闭环，下文给出一段基于系统结构化记录字段整理的代表性对话片段\footnote{对话片段为论文分析需要，对实验/仿真对话中的关键轮次进行脱敏、节选与行文规范化整理；其用途是呈现状态机推进与策略选择的因果链，不宣称逐字等同于某条原始运行日志。状态名与规则名的实现可在 \texttt{src/router.py} 与 \texttt{src/guardrails.py} 中定位。}。具体状态流转与对话如下：

\textbf{【系统干预记录编号：\texttt{Log-Buoyancy-012}】}
\begin{itemize}
    \item \textbf{学生画像}：P2（易受经验误导型）
    \item \textbf{主题}：浮力与深度
    \item \textbf{干预轮次}：5 轮
    \item \textbf{初始诊断迷思}：浮力随深度增大（压强经验投射）
    \item \textbf{系统策略序列}：\texttt{S3}（诊断）$\rightarrow$ \texttt{S4}（认知冲突）$\rightarrow$ \texttt{S5}（支架式引导）$\rightarrow$ \texttt{S6}（验证与加深）$\rightarrow$ \texttt{S7}（事实兜底/迁移检验）
    \item \textbf{干预结果}：学生完成自我解释，并进入 \texttt{S7} 的事实锚定与迁移检验阶段。
\end{itemize}

\begin{quote}
\textbf{[State: \texttt{S1 Guardrail\_Check} $\rightarrow$ \texttt{S3 Misconception\_Diagnosis}]} \\
\textbf{系统}：你提到“越深浮力越大”，能具体说说为什么这样想吗？ 

\textbf{[State: \texttt{S3} $\rightarrow$ \texttt{S4 Cognitive\_Conflict}]} \\
\textbf{学生}：因为水越深的地方，水的压力就越大啊，所以浮力肯定也跟着变大。\\
\textit{（诊断记录：单因归因错误置信度 0.92，策略匹配：\texttt{Consequence\_Exploration}）} \\
\textbf{系统}：按照这个想法，如果越深浮力越大，那么一个铁块在水里越往下沉，受到的浮力就越来越大。它越往下是不是越容易被这股越来越大的浮力给顶上来呢？ \\
\textbf{学生}：呃……好像不会被顶上来，铁块还是会一直沉到底。这里好像有点不对劲。

\textbf{[State: \texttt{S4} $\rightarrow$ \texttt{S5 Scaffolding\_Guidance}]} \\
\textbf{系统}：对呀，那为什么它没有被越来越大的浮力顶上来呢？我们可以控制变量比较一下：同一个铁块，完全没入水中后，在浅处和深处，它排开的水的体积变了吗？ \\
\textbf{学生}：没变，铁块的大小又没变。

\textbf{[State: \texttt{S5} $\rightarrow$ \texttt{S6 Verification\_Deepening}]} \\
\textbf{系统}：既然排开水的体积没变，那浮力会变吗？你觉得浮力到底和深度有关，还是和排开水的体积有关？

\textbf{[State: \texttt{S6} $\rightarrow$ \texttt{S7 Fact\_Grounding}]} \\
\textbf{学生}：哦！只要完全没入水中，排开的水体积就不变，所以浮力是不变的。浮力跟深度没关系，跟排开水的体积有关！\\
\textit{（诊断记录：理解验证通过，准许推进至 \texttt{S7}）} \\
\textbf{系统}：解释得非常好！那么如果是一艘潜水艇从水下 50 米潜到 100 米，它受到的浮力会变化吗？为什么？
\end{quote}

该代表性片段表明，归谬式后果探索之所以有效，在于它不是强行否定学生，而是让学生沿着自己的规则继续推演，并在推演中自行看到矛盾。对于稳定性较强的迷思概念，这种“让旧观念自己显出局限”的方式，往往比直接纠错更符合概念转变的认知机制\cite{posner1982accommodation,chi1994things}。

\subsection{案例二：教学支架豁免在缓解认知过载中的实际作用分析}
\label{sec:4.6.2}

第二个案例聚焦于系统在严格执行“不直接给答”原则时，如何通过有限支架豁免避免学生因持续卡顿而退出思考。以某一学生在浮力概念对话中的表现为例：该学生能够意识到“重的东西不一定沉”，但在进一步解释“浮力与什么有关”时反复停留在“和深度、和面积、和水的压力都有关系”的模糊表述中。若系统持续要求其自行说明，学生可能出现高负荷、低效循环，导致对话停滞。

在这一情境下，系统检测到学生已进行了多轮尝试，具有一定参与度，但仍无法组织出较清晰的解释，于是触发“确认性总结”与“思想实验”两种低强度支架豁免。具体状态流转与对话转录如下：

\textbf{【系统干预记录编号：\texttt{Log-Buoyancy-028}】}
\begin{itemize}
    \item \textbf{学生画像}：P3（高认知负荷/易挫败型）
    \item \textbf{主题}：浮力多因素混淆
    \item \textbf{干预轮次}：8 轮（节选后段）
    \item \textbf{初始诊断迷思}：多因素干扰导致解释混乱
    \item \textbf{系统策略序列}：\texttt{S4}（连续追问卡顿）$\rightarrow$ 触发 \texttt{ANTI\_LOOP\_RULES} 降级 $\rightarrow$ \texttt{S5}（思想实验支架）$\rightarrow$ \texttt{S6}（自我解释验证）
    \item \textbf{干预结果}：通过支架豁免成功化解认知死锁，恢复对话推进。
\end{itemize}

\begin{quote}
\textbf{[State: \texttt{S4 Cognitive\_Conflict} (Loop count: 2)]} \\
\textbf{系统}：那你觉得浮力跟什么有关呢？能举个例子说明吗？ \\
\textbf{学生}：应该和深度、和物体的面积、和水的压力都有关系吧……感觉好复杂，我不知道该怎么说。 \\
\textbf{系统}：那我们一个一个看，你觉得接触水面积大，浮力就一定大吗？ \\
\textbf{学生}：可能吧，我被搞晕了，完全想不清楚了。

\textbf{[State: 触发 \texttt{ANTI\_LOOP\_RULES} 规则 2，降级至 \texttt{S5 Scaffolding\_Guidance}]} \\
\textit{（日志诊断记录：情感分析=焦虑/挫败，触发思想实验支架豁免）} \\
\textbf{系统}：没关系，感觉有些复杂是很正常的。你刚才其实已经意识到，单看深度不够，还要看物体到底排开了多少水。（\textit{豁免动作：确认性总结}）我们来做一个简单的思想实验：如果把同一个小球放进两杯同样的水里，一杯浅一点，一杯深一点，只要小球都完全浸没，它排开的水量会变吗？（\textit{豁免动作：受控思想实验}） \\
\textbf{学生}：不会变，都是那个小球那么大。 \\
\textbf{系统}：既然排开的水量一样，那这两杯水里，小球受到的浮力一样吗？

\textbf{[State: \texttt{S5} $\rightarrow$ \texttt{S6 Verification\_Deepening}]} \\
\textbf{学生}：既然排开的水量一样，那浮力肯定一样。所以我之前想复杂了，其实只要看排开的水量就行了！
\end{quote}

如上所示，这种支架并没有直接替学生说完结论，而是降低了学生在变量控制、语言组织和工作记忆上的额外负担，使其得以把注意力重新集中到关键关系上。

从教学意义上看，支架豁免并不是对“不直接给答”原则的削弱，而是对其更精细的落实。真正需要避免的不是“系统说得稍多一点”，而是“系统替学生完成了关键推理”。当支架仅仅帮助学生减少无关负担、聚焦关键变量、组织已有想法时，它实际上增强了学生的元认知参与，而非削弱之。元认知辅导研究同样表明，有效的帮助不是简单增加信息量，而是优化学生求助与反思的方式\cite{aleven2004toward,roll2011improving}。因此，本研究中的支架豁免机制可以被视为在教学安全边界内对认知过载进行调节的一种必要设计。失败与边界案例（例如护栏连续触发仍无法恢复有效互动、或事实兜底后学生仍拒绝参与）将在第5章结合评测指标进一步讨论。
